\documentclass[]{article}
\usepackage{lmodern}
\usepackage{amssymb,amsmath}
\usepackage{ifxetex,ifluatex}
\usepackage{fixltx2e} % provides \textsubscript
\ifnum 0\ifxetex 1\fi\ifluatex 1\fi=0 % if pdftex
  \usepackage[T1]{fontenc}
  \usepackage[utf8]{inputenc}
\else % if luatex or xelatex
  \ifxetex
    \usepackage{mathspec}
  \else
    \usepackage{fontspec}
  \fi
  \defaultfontfeatures{Ligatures=TeX,Scale=MatchLowercase}
\fi
% use upquote if available, for straight quotes in verbatim environments
\IfFileExists{upquote.sty}{\usepackage{upquote}}{}
% use microtype if available
\IfFileExists{microtype.sty}{%
\usepackage{microtype}
\UseMicrotypeSet[protrusion]{basicmath} % disable protrusion for tt fonts
}{}
\usepackage[left=1.1in, right=1.1in, top=1.1in, bottom=1.1in]{geometry}
\usepackage{hyperref}
\hypersetup{unicode=true,
            pdftitle={Stat 435 Lecture Notes 4a},
            pdfauthor={Xiongzhi Chen; Washington State University},
            pdfborder={0 0 0},
            breaklinks=true}
\urlstyle{same}  % don't use monospace font for urls
\usepackage{color}
\usepackage{fancyvrb}
\newcommand{\VerbBar}{|}
\newcommand{\VERB}{\Verb[commandchars=\\\{\}]}
\DefineVerbatimEnvironment{Highlighting}{Verbatim}{commandchars=\\\{\}}
% Add ',fontsize=\small' for more characters per line
\usepackage{framed}
\definecolor{shadecolor}{RGB}{248,248,248}
\newenvironment{Shaded}{\begin{snugshade}}{\end{snugshade}}
\newcommand{\KeywordTok}[1]{\textcolor[rgb]{0.13,0.29,0.53}{\textbf{{#1}}}}
\newcommand{\DataTypeTok}[1]{\textcolor[rgb]{0.13,0.29,0.53}{{#1}}}
\newcommand{\DecValTok}[1]{\textcolor[rgb]{0.00,0.00,0.81}{{#1}}}
\newcommand{\BaseNTok}[1]{\textcolor[rgb]{0.00,0.00,0.81}{{#1}}}
\newcommand{\FloatTok}[1]{\textcolor[rgb]{0.00,0.00,0.81}{{#1}}}
\newcommand{\ConstantTok}[1]{\textcolor[rgb]{0.00,0.00,0.00}{{#1}}}
\newcommand{\CharTok}[1]{\textcolor[rgb]{0.31,0.60,0.02}{{#1}}}
\newcommand{\SpecialCharTok}[1]{\textcolor[rgb]{0.00,0.00,0.00}{{#1}}}
\newcommand{\StringTok}[1]{\textcolor[rgb]{0.31,0.60,0.02}{{#1}}}
\newcommand{\VerbatimStringTok}[1]{\textcolor[rgb]{0.31,0.60,0.02}{{#1}}}
\newcommand{\SpecialStringTok}[1]{\textcolor[rgb]{0.31,0.60,0.02}{{#1}}}
\newcommand{\ImportTok}[1]{{#1}}
\newcommand{\CommentTok}[1]{\textcolor[rgb]{0.56,0.35,0.01}{\textit{{#1}}}}
\newcommand{\DocumentationTok}[1]{\textcolor[rgb]{0.56,0.35,0.01}{\textbf{\textit{{#1}}}}}
\newcommand{\AnnotationTok}[1]{\textcolor[rgb]{0.56,0.35,0.01}{\textbf{\textit{{#1}}}}}
\newcommand{\CommentVarTok}[1]{\textcolor[rgb]{0.56,0.35,0.01}{\textbf{\textit{{#1}}}}}
\newcommand{\OtherTok}[1]{\textcolor[rgb]{0.56,0.35,0.01}{{#1}}}
\newcommand{\FunctionTok}[1]{\textcolor[rgb]{0.00,0.00,0.00}{{#1}}}
\newcommand{\VariableTok}[1]{\textcolor[rgb]{0.00,0.00,0.00}{{#1}}}
\newcommand{\ControlFlowTok}[1]{\textcolor[rgb]{0.13,0.29,0.53}{\textbf{{#1}}}}
\newcommand{\OperatorTok}[1]{\textcolor[rgb]{0.81,0.36,0.00}{\textbf{{#1}}}}
\newcommand{\BuiltInTok}[1]{{#1}}
\newcommand{\ExtensionTok}[1]{{#1}}
\newcommand{\PreprocessorTok}[1]{\textcolor[rgb]{0.56,0.35,0.01}{\textit{{#1}}}}
\newcommand{\AttributeTok}[1]{\textcolor[rgb]{0.77,0.63,0.00}{{#1}}}
\newcommand{\RegionMarkerTok}[1]{{#1}}
\newcommand{\InformationTok}[1]{\textcolor[rgb]{0.56,0.35,0.01}{\textbf{\textit{{#1}}}}}
\newcommand{\WarningTok}[1]{\textcolor[rgb]{0.56,0.35,0.01}{\textbf{\textit{{#1}}}}}
\newcommand{\AlertTok}[1]{\textcolor[rgb]{0.94,0.16,0.16}{{#1}}}
\newcommand{\ErrorTok}[1]{\textcolor[rgb]{0.64,0.00,0.00}{\textbf{{#1}}}}
\newcommand{\NormalTok}[1]{{#1}}
\usepackage{graphicx,grffile}
\makeatletter
\def\maxwidth{\ifdim\Gin@nat@width>\linewidth\linewidth\else\Gin@nat@width\fi}
\def\maxheight{\ifdim\Gin@nat@height>\textheight\textheight\else\Gin@nat@height\fi}
\makeatother
% Scale images if necessary, so that they will not overflow the page
% margins by default, and it is still possible to overwrite the defaults
% using explicit options in \includegraphics[width, height, ...]{}
\setkeys{Gin}{width=\maxwidth,height=\maxheight,keepaspectratio}
\IfFileExists{parskip.sty}{%
\usepackage{parskip}
}{% else
\setlength{\parindent}{0pt}
\setlength{\parskip}{6pt plus 2pt minus 1pt}
}
\setlength{\emergencystretch}{3em}  % prevent overfull lines
\providecommand{\tightlist}{%
  \setlength{\itemsep}{0pt}\setlength{\parskip}{0pt}}
\setcounter{secnumdepth}{0}
% Redefines (sub)paragraphs to behave more like sections
\ifx\paragraph\undefined\else
\let\oldparagraph\paragraph
\renewcommand{\paragraph}[1]{\oldparagraph{#1}\mbox{}}
\fi
\ifx\subparagraph\undefined\else
\let\oldsubparagraph\subparagraph
\renewcommand{\subparagraph}[1]{\oldsubparagraph{#1}\mbox{}}
\fi

%%% Use protect on footnotes to avoid problems with footnotes in titles
\let\rmarkdownfootnote\footnote%
\def\footnote{\protect\rmarkdownfootnote}

%%% Change title format to be more compact
\usepackage{titling}

% Create subtitle command for use in maketitle
\newcommand{\subtitle}[1]{
  \posttitle{
    \begin{center}\large#1\end{center}
    }
}

\setlength{\droptitle}{-2em}

  \title{Stat 435 Lecture Notes 4a}
    \pretitle{\vspace{\droptitle}\centering\huge}
  \posttitle{\par}
    \author{Xiongzhi Chen \\ Washington State University}
    \preauthor{\centering\large\emph}
  \postauthor{\par}
    \date{}
    \predate{}\postdate{}
  
\usepackage{bbm}
\usepackage{amssymb}
\usepackage{amsmath}
\usepackage{graphicx,float}
\usepackage{hyperref}

\begin{document}
\maketitle

{
\setcounter{tocdepth}{2}
\tableofcontents
}
\section{}\label{section}

\section{Model assessment}\label{model-assessment}

\subsection{Model and estimate}\label{model-and-estimate}

\begin{itemize}
\tightlist
\item
  Model:
  \(Y=\beta_0+\beta_1 X_1 + \beta_2 X_2 + \ldots + \beta_p X_p + \varepsilon\)
\item
  Observations:
  \(\mathbf{z}_i=(y_i,x_{i1},x_{i2},\ldots,x_{ip}),i=1,\ldots,n\), where
  \(y_i\) is the \(i\)th observation for \(Y\) and \(x_{ij}\) that for
  \(X_j\)
\item
  Estimate
  \(\hat{\boldsymbol{\beta}}=(\hat{\beta}_0,\hat{\beta}_1,\ldots,\hat{\beta}_p)\)
  of \({\boldsymbol{\beta}}=({\beta}_0,{\beta}_1,\ldots,{\beta}_p)\)
\item
  Fitted model:
  \(\hat{y}_i=\hat{\beta}_0+\hat{\beta}_1 x_{i1} + \hat{\beta}_2 x_{i2} + \ldots + \hat{\beta}_p x_{ip}\)
\item
  Residuals: \(e_i=y_i - \hat{y}_i\)
\end{itemize}

\emph{Note:} there are many ways to obtain an estimate
\(\hat{\boldsymbol{\beta}}\) of \({\boldsymbol{\beta}}\)

\subsection{Training error}\label{training-error}

\begin{itemize}
\item
  \emph{Training error} is a measure on the performance of a model when
  it is fitted/trained from a set of observations that we call the
  \emph{training set}. Namely, \emph{training error measures how well a
  fitted/trained model fits/learns the training set}.
\item
  A commonly used training error is the \emph{residual sum of squares
  (RSS)} (though other choices are available)
\item
  For example, if we take the set of \(n\) observations as the training
  set and use RSS as the training error, then the \emph{least squares
  estimate (LSE)} minimizes \(\textrm{RSS}=\sum_{i=1}^m e_i^2\)
\end{itemize}

\subsection{Test error}\label{test-error}

\begin{itemize}
\tightlist
\item
  \emph{Test error} is a measure on the \emph{predictive performance} of
  a (fitted/trained) model when it is used to predict the responses of
  new observations on the predictors \(\tilde{X}=(X_1,\ldots,X_p)\) that
  are not in the training set (that is used to fit/train the model).
  Namely, \emph{test error measures how well a fitted/trained model
  predicts responses for observations that are not part of a training
  set}.
\item
  For example, given a new observation
  \(\mathbf{x}_0=(x_{01},x_{02},\ldots,x_{0p})\) for \(\tilde{X}\), let
  \(y_0\) be the predicted response given by the fitted/trained model.
  Then we can use \(E(y_0 -\hat{y}_0)^2\) as the test error.
\end{itemize}

\emph{Note:} other measures of test error are available

\subsection{Training error and test
error}\label{training-error-and-test-error}

\begin{itemize}
\tightlist
\item
  Training error measures how well a fitted/trained model fits/learns
  the \emph{seen} observations (of the training set), whereas test error
  measures how well this model \emph{predicts unseen responses} based on
  unseen observations on predictors
\item
  \emph{Test error is alway unknown and needs to be estimated}
\item
  A model that has good predictive performance is referred to as having
  good ``generalization performance'' or ``generalizes well''
\item
  A model that fits the training set very well or perfectly is often
  referred to as ``overfitting'' or ``overfits''
\end{itemize}

\subsection{Training error and test
error}\label{training-error-and-test-error-1}

\begin{itemize}
\tightlist
\item
  Classic wisdom: it is believed that an overfitting model that fits the
  training set very well or perfectly (i.e., with very small training
  error) cannot generalize well (i.e., cannot have small testing error)
\item
  Modern wisdom: when there are relatively few parameters in the model,
  the classic wisdom is sensible. However, beyond a critical setting, an
  overfitting model (with many parameters) can have very good
  generalization performance.
\end{itemize}

\emph{Note:} Modern wisdom, dubbed as the ``\textbf{double descent
curve}'', was discovered by Dr.~Belkin and his coauthors

\subsection{Double descent curve}\label{double-descent-curve}

\begin{center}\includegraphics[width=0.85\linewidth]{doubledescent} \end{center}

\emph{Image credit:} Belkin et al; doi.org/10.1073/pnas.1903070116

\section{Cross-validation}\label{cross-validation}

\subsection{Overview}\label{overview}

\emph{Cross-validation is a resampling technique to estimate test
error}, and is often implemented as follows:

\begin{itemize}
\tightlist
\item
  Randomly divides a set of observations into a \emph{training set} and
  \emph{validation set}
\item
  Use the training set to fit a model (that optimizes some
  \emph{training error}), and apply the fitted model to predict the
  responses for the observations in the validation set to obtain an
  estimate of the \emph{test error} of the model
\item
  Do the above independently for different random splits, and estimate
  test error from the estimates of test error
\end{itemize}

\subsection{Pictorial description}\label{pictorial-description}

\begin{center}\includegraphics[width=0.9\linewidth]{fig5_1} \end{center}

\subsection{CV: estimate test error}\label{cv-estimate-test-error}

With \(n\) observations \(\mathbf{z}_i = (y_i,\mathbf{x}_i)\),

\begin{itemize}
\item
  Randomly split the \(n\) observations into a training set
  \(\mathcal{T}_1\) with \(n_1\) observations, and a validation set
  \(\mathcal{V}_1\) with \(n_2=n-n_1\) observations
\item
  Fit model \(M_l\) using \(\mathcal{T}\), apply fitted model
  \(\hat{M}_l\) to predict responses in \(\mathcal{V}_1\), and compute
  the \emph{mean squared error (MSE)}
  \[\textrm{MSE}(\mathcal{V}_1)=n_2^{-1} \sum_{y_i \in \mathcal{V}_1} (y_i -\hat{y}_i)^2,\]
  where \(\hat{y}_i\) is the fitted value for \(y_i\)
\end{itemize}

\emph{Note:} \(\mathcal{T}_1\) and \(\mathcal{V}_1\) are disjoint

\subsection{CV: estimate test error}\label{cv-estimate-test-error-1}

\begin{itemize}
\tightlist
\item
  Repeat the above ``splitting-training-validating'' steps independently
  \(k\) times for model \(M_l\) to obtain \(k\) MSE's
  \(\textrm{MSE}(\mathcal{V}_j),j=1,\ldots,k\), and estimate the test
  \emph{mean squared error (MSE)} of model \(M_l\) by
  \[\textrm{CV}(M_l)=k^{-1} \sum_{j=1}^k \textrm{MSE}(\mathcal{V}_j)\]
\end{itemize}

\emph{Note:} the above steps give an estimate of test error of model
\(M_l\)

\subsection{CV: model selection}\label{cv-model-selection}

\begin{itemize}
\tightlist
\item
  Obtain via cross-validation the estimated test error for each model
  \(M_l,l=1,\ldots,L\)
\item
  Pick the model that has the smallest estimated test MSE, i.e., pick
  the optimal model \(M^{\ast}\) such that
  \[\textrm{CV}(M^*)=\min_{1\le l \le L}\textrm{CV}(M_l)\]
\end{itemize}

\emph{Note:} the above gives the best model among a set of models

\subsection{Illustration}\label{illustration}

\begin{center}\includegraphics[width=0.9\linewidth]{fig5_2} \end{center}

\subsection{\texorpdfstring{\(k\)-fold
cross-validation}{k-fold cross-validation}}\label{k-fold-cross-validation}

The \(k\)-fold CV for model \(M_l\) is implemented as follows:

\begin{itemize}
\tightlist
\item
  Randomly splits the data set into \(k\) groups, or ``folds'', of
  approximately equal sizes
\item
  Pick a fold, say, folder \(j\), as the validation set (i.e., the
  ``held-out'' fold), fit the model on the remaining \(k-1\) folds, and
  compute the mean squared error, \(\textrm{MSE}_{j}\), on the
  observations in the held-out fold
\item
  Do the above for all \(j,j=1,\ldots,k\), and estimate the test MSE of
  the model by \[
  \textrm{CV}({M_l})=\frac{1}{k}\sum_{j=1}^{k}\textrm{MSE}_{j}
  \]
\end{itemize}

\section{Variable and model
selection}\label{variable-and-model-selection}

\subsection{Motivation}\label{motivation}

Settings:

\begin{itemize}
\item
  Model:
  \(Y=\beta_0+\beta_1 X_1 + \beta_2 X_2 + \ldots + \beta_p X_p + \varepsilon\)
\item
  Observations: \((y_i,x_{i1},x_{i2},\ldots,x_{ip}),i=1,\ldots,n\),
  where \(y_i\) is the \(i\)th observation for \(Y\) and \(x_{ij}\) that
  for \(X_j\)
\item
  Estimate
  \(\hat{\boldsymbol{\beta}}=(\hat{\beta}_0,\hat{\beta}_1,\ldots,\hat{\beta}_p)\)
  of \({\boldsymbol{\beta}}=({\beta}_0,{\beta}_1,\ldots,{\beta}_p)\)
\end{itemize}

\emph{How to obtain an estimate \(\hat{\boldsymbol{\beta}}\) of
\({\boldsymbol{\beta}}\) depends critically on}

\begin{itemize}
\tightlist
\item
  the relative magnitudes of \(p+1\) and \(n\)
\item
  what properties we want \(\hat{\boldsymbol{\beta}}\) to have
\end{itemize}

\subsection{\texorpdfstring{Classic scenario:
\(p+1 \le n\)}{Classic scenario: p+1 \textbackslash{}le n}}\label{classic-scenario-p1-le-n}

When the number of parameters is not larger than the sample size:

\begin{itemize}
\tightlist
\item
  The \emph{least squares estimate (LSE)} is \emph{uniquely} defined
\item
  The LSE is unbiased, i.e.,
  \(E(\hat{\boldsymbol{\beta}})={\boldsymbol{\beta}}\), when
  \(E(\varepsilon)=0\)
\item
  The LSE is optimal in some sense (e.g., Gauss-Markov theorem, meaning
  that the LSE has the smallest variance among all linear unbiased
  estimators, if the random errors are uncorrelated and have equal
  variances and zero expectation)
\end{itemize}

\subsection{\texorpdfstring{Classic scenario:
\(p+1 \le n\)}{Classic scenario: p+1 \textbackslash{}le n}}\label{classic-scenario-p1-le-n-1}

When the number of parameters is not larger than the sample size but
\emph{there are many potential predictors}, we often desire a
\emph{small} model that is easy to interpret and perform well. Namely,
we still need to consider:

\begin{itemize}
\tightlist
\item
  Which predictors are the \emph{most relevant} to the response
\item
  how to build a \emph{small} model using a few predictors that can well
  predict the response
\end{itemize}

\emph{Variable or model selection is needed in the classic scenario when
there are many potential predictors.}

\subsection{\texorpdfstring{Modern scenario:
\(p+1 > n\)}{Modern scenario: p+1 \textgreater{} n}}\label{modern-scenario-p1-n}

When the number of parameters is bigger than the sample size,

\begin{itemize}
\tightlist
\item
  The LSE is \emph{not uniquely} defined, i.e., there are infinitely
  many models that minimizes the residual sum of squares
\item
  An LSE is biased in general, i.e.,
  \(E(\hat{\boldsymbol{\beta}}) \ne {\boldsymbol{\beta}}\), even when
  \(E(\varepsilon)=0\)
\item
  An LSE has infinite variance (and Gauss-Markov theorem does not hold
  for LSE)
\end{itemize}

\subsection{\texorpdfstring{Modern scenario:
\(p+1 > n\)}{Modern scenario: p+1 \textgreater{} n}}\label{modern-scenario-p1-n-1}

In this scenario, we have a few choices:

\begin{itemize}
\tightlist
\item
  perform variable/model selection: select a few most relevant
  predictors to form a model and then apply LSE to estimate the model
  parameters
\item
  employ different estimation methods: use a different method (such as
  \emph{shrinkage}) than LSE to estimate coefficients of predictors in a
  model (as a \emph{bias-variance trade-off})
\item
  use dimension reduction: projecting all predictors onto a smaller
  subspace and use transformed predictors to build a model
\end{itemize}

\section{Best subset selection}\label{best-subset-selection}

\subsection{Overview}\label{overview-1}

\begin{itemize}
\item
  Consider a linear model with \(p\) predictors: \[
  Y=\beta_0+\beta_1 X_1 + \beta_2 X_2 + \ldots + \beta_p X_p + \varepsilon
  \]
\item
  \emph{Best subset selection} is a ``brute-force'' method that checks
  \emph{each of the \(2^p\) possible linear submodels} and picks the
  best one under some criterion
\item
  Best subset selection indeed gives the \emph{best subset of predictors
  (in terms of linear model)} among all \(p\) predictors under a
  criterion
\end{itemize}

\subsection{Criteria for subset
selection}\label{criteria-for-subset-selection}

Some criteria for subset/variable selection:

\begin{itemize}
\tightlist
\item
  Mallow's \(C_p\)
\item
  AIC (i.e., Akaike information criterion),
\item
  BIC (i.e., Bayesian information criterion)
\item
  Adjusted \(R^2\)
\item
  Cross-validated prediction error
\end{itemize}

\subsection{Implementation}\label{implementation}

\begin{center}\includegraphics[width=0.85\linewidth]{Alg6.1} \end{center}

\subsection{Illustration}\label{illustration-1}

\begin{center}\includegraphics[width=0.9\linewidth]{fig6.1} \end{center}

\section{Forward/backwards stepwise
selection}\label{forwardbackwards-stepwise-selection}

\subsection{Overview}\label{overview-2}

Forward stepwise selection

\begin{itemize}
\tightlist
\item
  only fits
  \[1+\sum_{k=0}^{p-1}\left(p-k\right)=1+2^{-1}p\left(p+1\right)\]
  \emph{linear submodels} out of a total of \(2^p\) linear submodels
\item
  is not guaranteed to find the best model among all linear submodels
\item
  has much smaller computational intensity than best subset selection
\end{itemize}

\subsection{Implementation}\label{implementation-1}

\begin{center}\includegraphics[width=0.9\linewidth]{alg6_2} \end{center}

\subsection{Illustration}\label{illustration-2}

\begin{center}\includegraphics[width=0.9\linewidth]{t6_1} \end{center}

\subsection{Backwards stepwise
selection}\label{backwards-stepwise-selection}

\begin{center}\includegraphics[width=0.9\linewidth]{alg6_3} \end{center}

\section{Choosing the optimal model}\label{choosing-the-optimal-model}

\subsection{General principle}\label{general-principle}

For best subset, forward stepwise, and backwards stepwise selections,
\emph{we need to select a best model from the best submodels}. However,
neither \emph{training} set RSS nor \emph{training} set \(R^{2}\) can be
used for this purpose since

\begin{itemize}
\tightlist
\item
  neither ensures good predictive performance in terms of test error of
  the resultant model
\item
  either tends to pick a model that has the largest possible size in
  terms of the number of parameters
\end{itemize}

So, a practical way to select a single best model is to \textbf{balance
the RSS on the training error and the model size}

\subsection{Four methods}\label{four-methods}

Four methods to choose the optimal model:

\begin{itemize}
\tightlist
\item
  Mallow's \(C_p\)
\item
  Akaike information criterion
\item
  Bayesian information criterion
\item
  Cross-validation
\end{itemize}

\subsection{\texorpdfstring{Mallow's
\(C_p\)}{Mallow's C\_p}}\label{mallows-c_p}

Mallow's

\begin{equation}
C_{p}=\frac{1}{n}\left( \textrm{RSS}+2d\hat{\sigma}^{2}\right)
\end{equation}

\begin{itemize}
\tightlist
\item
  \(d\) is the number of predictors in the model
\item
  Typically \(\hat{\sigma}^{2}\), an estimate of
  \(\sigma^2=\textrm{Var}(\varepsilon)\), is estimated using the full
  model containing all predictors
\end{itemize}

\emph{Remark:} If \(\hat{\sigma}^{2}\) is an unbiased estimate of
\(\sigma^{2}\), then \(C_{p}\) is an unbiased estimate of test MSE

\subsection{Akaike information
criterion}\label{akaike-information-criterion}

\begin{itemize}
\tightlist
\item
  The \emph{Akaike information criterion (AIC)} is mainly used together
  with the \emph{maximum likelihood method}
\item
  For a linear model with Gaussian errors, the \emph{maximum likelihood
  estimate (MLE)} is the same as the \emph{least squares estimate
  (LSE)}. In this case, AIC is, up to an additive constant, given by \[
  \textrm{AIC}=\frac{1}{n\hat{\sigma}^{2}}\left(  \textrm{RSS}+2d\hat{\sigma}^{2}\right)
  \]
\end{itemize}

\subsection{Bayesian information
criterion}\label{bayesian-information-criterion}

\begin{itemize}
\tightlist
\item
  The \emph{Bayesian information criterion (BIC)} is derived from a
  Bayesian perspective
\item
  For the least squares model with \(d\) parameters, the BIC is, up to
  irrelevant constants, given by \[
  \textrm{BIC}=\frac{1}{n\hat{\sigma}^{2}}\left( \textrm{RSS}+d\hat{\sigma}^{2}\log n\right)
  \]
\item
  BIC generally places a heavier penalty on models with many variables.
  For large \(n\), BIC is bigger than \(C_{p}\)
\end{itemize}

\subsection{\texorpdfstring{Adjusted
\(R^{2}\)}{Adjusted R\^{}\{2\}}}\label{adjusted-r2}

\begin{itemize}
\tightlist
\item
  The adjusted \(R^{2}\) is given by \[
  \text{Adjusted }R^{2}=1-\frac{\textrm{RSS}/\left(  n-d-1\right)  }{\textrm{TSS}/\left(n-1\right)  }
  \]
\item
  \emph{Unlike \(C_{p}\), AIC and BIC, for which a small value indicates
  a model with a low test error, a large value of adjusted \(R^{2}\)
  indicates a model with a low test error}
\end{itemize}

The intuition behind the adjusted \(R^{2}\) is that ``once all of the
correct variables have been included in the model, adding additional
noise variables will lead to a very small decrease in RSS''

\subsection{Recap}\label{recap}

Let \(d\) be number of predictors in model, \(n\) sample size, and
\(\hat{\sigma}^{2}\) an estimate of
\(\sigma^2=\textrm{Var}(\varepsilon)\):

\begin{itemize}
\item
  \(C_{p}=\frac{1}{n}\left( \textrm{RSS}+2d\hat{\sigma}^{2}\right)\)
\item
  \(\textrm{AIC}=\frac{1}{n\hat{\sigma}^{2}}\left( RSS+2d\hat{\sigma}^{2}\right)\)
\item
  \(\textrm{BIC}=\frac{1}{n\hat{\sigma}^{2}}\left( \textrm{RSS}+d\hat{\sigma}^{2}\log n\right)\)
\item
  \(\text{Adjusted }R^{2}=1-\frac{\textrm{RSS}/\left( n-d-1\right) }{\textrm{TSS}/\left(n-1\right) }\)
\end{itemize}

All formulae for \(C_{p}\), AIC and BIC are \emph{for a linear model fit
using least squares}; \(C_p\), AIC and BIC all have good theoretical
justifications

\subsection{Illustration}\label{illustration-3}

\begin{center}\includegraphics[width=0.9\linewidth]{fig6_2} \end{center}

\subsection{\texorpdfstring{\(k\)-fold
cross-validation}{k-fold cross-validation}}\label{k-fold-cross-validation-1}

Recall CV for model selection:

\begin{itemize}
\tightlist
\item
  Obtain via cross-validation the estimated test error for each model
  \(M_l,l=1,\ldots,L\)
\item
  Pick the model that has the smallest estimated test MSE, i.e., pick
  the optimal model \(M^{\ast}\) such that
  \[\textrm{CV}(M^*)=\min_{1\le l \le L}\textrm{CV}(M_l)\]
\end{itemize}

\subsection{\texorpdfstring{\(k\)-fold
cross-validation}{k-fold cross-validation}}\label{k-fold-cross-validation-2}

The \(k\)-fold CV for model \(M_l\) is implemented as follows:

\begin{itemize}
\tightlist
\item
  Randomly splits the data set into \(k\) groups, or ``folds'', of
  approximately equal sizes
\item
  Pick a fold, say, folder \(j\), as the validation set (i.e., the
  ``held-out'' fold), fit the model on the remaining \(k-1\) folds, and
  compute the mean squared error, \(\textrm{MSE}_{j}\), on the
  observations in the held-out fold
\item
  Do the above for all \(j,j=1,\ldots,k\), and estimate the test MSE of
  the model by \[
  \textrm{CV}({M_l})=\frac{1}{k}\sum_{j=1}^{k}\textrm{MSE}_{j}
  \]
\end{itemize}

\subsection{\texorpdfstring{\(k\)-fold
cross-validation}{k-fold cross-validation}}\label{k-fold-cross-validation-3}

Guideline:

\begin{itemize}
\tightlist
\item
  Usually \(k=5\) or \(10\) is chosen, as these values have been shown
  empirically to yield test error rate estimates that suffer neither
  from excessively high bias nor from very high variance
\end{itemize}

\subsection{\texorpdfstring{\(k\)-fold
cross-validation}{k-fold cross-validation}}\label{k-fold-cross-validation-4}

\begin{center}\includegraphics[width=0.9\linewidth]{fig6_3} \end{center}

\subsection{One-standard-error rule}\label{one-standard-error-rule}

Different random splitting schemes often lead to differential optimal
models. So,

\begin{itemize}
\item
  first, we can calculate the standard error of the estimated test MSE
  for each model size, by repeatedly validating ``the best model'' of
  this model size
\item
  then select the smallest model for which the estimated test error is
  within one standard error of the lowest point on the front curve for
  the estimated MSEs of ``the best models''
\end{itemize}

\subsection{License and session
Information}\label{license-and-session-information}

\href{http://math.wsu.edu/faculty/xchen/stat412/LICENSE.html}{License}

\begin{Shaded}
\begin{Highlighting}[]
\NormalTok{>}\StringTok{ }\KeywordTok{sessionInfo}\NormalTok{()}
\NormalTok{R version }\FloatTok{3.5.0} \NormalTok{(}\DecValTok{2018-04-23}\NormalTok{)}
\NormalTok{Platform:}\StringTok{ }\NormalTok{x86_64-w64-mingw32/}\KeywordTok{x64} \NormalTok{(}\DecValTok{64}\NormalTok{-bit)}
\NormalTok{Running under:}\StringTok{ }\NormalTok{Windows }\DecValTok{10} \KeywordTok{x64} \NormalTok{(build }\DecValTok{19043}\NormalTok{)}

\NormalTok{Matrix products:}\StringTok{ }\NormalTok{default}

\NormalTok{locale:}
\NormalTok{[}\DecValTok{1}\NormalTok{] LC_COLLATE=English_United States}\FloatTok{.1252} 
\NormalTok{[}\DecValTok{2}\NormalTok{] LC_CTYPE=English_United States}\FloatTok{.1252}   
\NormalTok{[}\DecValTok{3}\NormalTok{] LC_MONETARY=English_United States}\FloatTok{.1252}
\NormalTok{[}\DecValTok{4}\NormalTok{] LC_NUMERIC=C                          }
\NormalTok{[}\DecValTok{5}\NormalTok{] LC_TIME=English_United States}\FloatTok{.1252}    

\NormalTok{attached base packages:}
\NormalTok{[}\DecValTok{1}\NormalTok{] stats     graphics  grDevices utils     datasets  methods  }
\NormalTok{[}\DecValTok{7}\NormalTok{] base     }

\NormalTok{other attached packages:}
\NormalTok{[}\DecValTok{1}\NormalTok{] knitr_1}\FloatTok{.21}

\NormalTok{loaded via a }\KeywordTok{namespace} \NormalTok{(and not attached):}
\StringTok{ }\NormalTok{[}\DecValTok{1}\NormalTok{] compiler_3}\FloatTok{.5.0}  \NormalTok{magrittr_1}\FloatTok{.5}    \NormalTok{tools_3}\FloatTok{.5.0}    
 \NormalTok{[}\DecValTok{4}\NormalTok{] htmltools_0}\FloatTok{.3.6} \NormalTok{yaml_2}\FloatTok{.2.0}      \NormalTok{Rcpp_1}\FloatTok{.0.3}     
 \NormalTok{[}\DecValTok{7}\NormalTok{] stringi_1}\FloatTok{.2.4}   \NormalTok{rmarkdown_1}\FloatTok{.11}  \NormalTok{stringr_1}\FloatTok{.3.1}  
\NormalTok{[}\DecValTok{10}\NormalTok{] xfun_0}\FloatTok{.4}        \NormalTok{digest_0}\FloatTok{.6.18}   \NormalTok{evaluate_0}\FloatTok{.12}  
\end{Highlighting}
\end{Shaded}


\end{document}
